

Algorithmic differentiation  is a class of computational methods
which leverage the chain rule in order to exactly evaluate\footnote{Within
  machine precision.} the gradient of a mathematical computation which is
implemented as a computer program. 
These methods generally leverage either (1) source code transformation
techniques or (2) operator overloading so that such a gradient
computation can be deduced entirely from the natural semantics of the
programming language being used. 
Much of the early literature which explores applications of \autodiff{} to fields like structural modeling
\autocite{ozaki1995higherorder,martins2001connection} employ the first
approach, where prior to compilation, an pre-processing step is
performed on some source code which generates new code implementing the
derivative of an operation. This approach is generally used with
primitive languages like FORTRAN and C which lack built-in abstractions
like operator overloading and was likely selected for these works
because at the time, such languages were the standard in the field.
Although these compile-time methods allow for powerful optimizations to
be used, they can add significant complexity to a project. 
Languages like C++, Python and Matlab allow user-defined types to \emph{overload}
built-in primitive operators like \texttt{+} and \texttt{*} so that when
they are applied to certain data types they invoke a user-specified
procedure. 
These features allow for \autodiff{} implementations which are
entirely contained within a programming language's standard ecosystem.

\section{Literature}
These methods are often further
categorized into \emph{forward} and \emph{reverse} modes, the later of
which further specializes to the \emph{backpropagation} class of
algorithms which has become ubiquitous in the field of machine learning.
% Hybrid Differentiation in Finite Element Analysis\label{hybrid-differentiation-in-finite-element-analysis}

\subsection{Direct Differentiation in
Computational Mechanics}

The direct differentiation method (DDM) computes sensitivities
by analytically differentiating the finite element governing equations
with respect to model parameters. 
In the earthquake engineering and
reliability community, DDM has been a preferred approach for obtaining
exact response gradients. 
\cite{haukaas2003finite} implemented a comprehensive DDM-based sensitivity module in the \OpenSees{} platform, enabling accurate and efficient computation of
response sensitivities for material, load, and shape parameters.

In such implementations, one derives the partial derivative of each
element's stiffness matrix, internal force, etc., and assembles the
global gradient. 
The benefit is that the gradients are exact and often very efficient to evaluate. 

\cite{scott2004response}, 
applied DDM to nonlinear beam-column elements.

\cite{al-aukaily2017response}




\subsection{Automatic Differentiation in Computational Mechanics}

In the past decade, automatic differentiation (AD) has made
significant inroads into computational mechanics, offering a way to
obtain exact gradients \emph{without} manual derivation. 
\Autodiff{} treats the finite element
algorithm as a composition of differentiable operations and applies the
chain rule systematically. 

Historic accounts of the development
of \autodiff{} can be found in \textcite{iri2009automatic}. 
%
Such treatments generally begin with the work of Wengert (e.g. \textcite{wengert1964simple}). 
A more recent work,
\textcite{elliott2009beautiful} develops an elegant presentation and
implementation of forward-mode AD in the context of purely functional
programming.

A thorough treatment of \autodiff{} for iterative procedures such as the use of
the Newton-Raphson method for solving nonlinear equations is presented
by \textcite{beck1994automatic}, \textcite{gilbert1992automatic} and
\textcite{griewank1993derivative}. 

Modern \autodiff{} tools can compute Jacobians and
gradients programmatically, and researchers have started
leveraging these tools in finite element contexts:

\begin{itemize}
\item
  \textbf{AD in FE Frameworks:} Large-scale multiphysics codes have
  begun to incorporate \autodiff{} to form exact system Jacobians. 
  A notable
  example is the MOOSE framework \citep{MOOSE} which uses the MetaPhysicL \autodiff{}
  library. 
  \cite{lindsay2021automatic} describe the integration of \autodiff{}
  into MOOSE's finite element routines for the automatic
  computation of consistent tangent matrices for nonlinear solves. 
  They report that this approach yields an exact Jacobian at relatively small overhead, enhancing the robustness and convergence of Newton--Raphson solvers. 
\item
  \textbf{Differentiable FE Libraries:} 
  \cite{liang2023pytorchfea} introduced
  \emph{PyTorch-FEA}, which implements finite element analysis using
  PyTorch tensor operations. 
  By ``taking advantage of autograd, an
  automatic differentiation mechanism in PyTorch,'' the library can
  seamlessly compute gradients of FE response with respect to input
  parameters. 
  This enabled them to tackle inverse problems by coupling
  the FE model with neural networks. 

  
  JAX-FEM \citep{xue2023jaxfem}  is a differentiable 3D FE solver built on Google's \Jax library \citep{frostig2018compiling}. 
  % It handled a 3D problem with 7.7 million
  % degrees of freedom and achieved an order-of-magnitude speedup over a
  % commercial FE code by running on GPU. 
  JAX-FEM uses reverse-mode \autodiff{} to
  obtain gradients, so one can perform automatic inverse design
  -- e.g.~topology optimization -- ``in a fully automatic manner
  without the need to manually derive sensitivities''. 
  % This
  % shows that with modern hardware and \autodiff{} tools, even large-scale FE
  % models can be made differentiable, opening the door to integrating FEM
  % with AI-driven design. 
  % (That said, JAX-FEM's success also relies on
  % problem structure and substantial computing power; fully
  % differentiating a million-DOF nonlinear problem is still extremely
  % memory-intensive in general.)
\item
  \textbf{AD in Specialized Methods:} Researchers have also applied \autodiff{}
  within specific computational mechanics techniques to simplify their
  implementation. 
  \cite{oberbichler2021efficient} used \autodiff{} to derive the internal force vector and consistent stiffness matrix for
  nonlinear isogeometric elements from a single energy
  functional. 
  They compared forward-mode vs.~reverse-mode
  \autodiff{} for this task and found the reverse (adjoint) mode far more
  efficient for high-degree elements: the reverse-mode Jacobian required
  significantly fewer operations and intermediate storage, and
  crucially, the number of AD intermediate variables no longer grew with
  the polynomial degree of shape functions. 
  In FE problems
  with many state variables but relatively fewer output quantities
  (e.g.~objective function or a few constraints), adjoint/reverse mode
  is usually optimal.
\end{itemize}

While \autodiff{} offers automation and exactness, it can
become computationally prohibitive for entire nonlinear finite element
simulations. 
The memory footprint of reverse-mode \autodiff{} is a known concern. 
In practice, fully
differentiable finite element codes often require careful engineering to be feasible. 
This is where hybrid approaches enter the picture -- combining analytic and
automatic differentiation to balance performance with flexibility.

\subsection{Hybrid Approaches}

Given the pros and cons of the two methods, a hybrid
differentiation strategy is attractive: use analytic
differentiation for parts of the problem where we can efficiently 
derive formulas, and use \emph{automatic differentiation} to handle the remainder. 
This idea has been explored in various forms.

% \textbf{Concept of Hybrid Differentiation:} 
\cite{bartlett2008hybrid} were early proponents of this approach in the context of
simulation codes. 
They observed that ``present AD tools\ldots rarely
compare to analytically derived versions'' in efficiency, but \autodiff{} is
far easier to implement. 
Thus, \emph{the focus of their work
was to combine the strength of these methods into a hybrid
strategy''}. 
In other words, one should ``select the
appropriate components'' of the simulation to differentiate manually
vs.~automatically to achieve an optimal balance of developer effort
and runtime cost. 
In their demonstration (fluid dynamics simulation in
C++), they hand-differentiated critical low-level routines and used
template-based AD for the rest -- yielding near-analytic performance
with much less manual coding.

% \textbf{Contemporary Examples:} 
Recently, \cite{dummer2024robust} put a hybrid approach into practice for nonlinear material modeling. 
They introduced a \emph{``semi-analytical split approach''}
  for a finite strain generalized continuum model. 
  In their framework, the global finite element equations were derived
  analytically as usual, but the material's stress-update
  algorithm was differentiated using automatic differentiation
  with hyper-dual numbers. 
  Hyper-dual \autodiff{} provided first and
  second derivatives of the stress update,
  yielding an exact consistent tangent matrix for the material
  sub-problem. 
  This split saved them from manually
  differentiating an arduous constitutive law, while still embedding
  that law into an efficient FE solver. 
  In effect, Dummer \emph{et al.} achieved a modular differentiation: the element
  stiffness contributions from a complicated material are obtained by 
  \autodiff{}, and everything else (element formulation, assembly, solver) uses
  standard analytic differentiation.

% \textbf{Adjoint-AD Hybrids:} 
Another approach to hybridization is mixing \emph{adjoint methods} with \autodiff{}. 
Adjoint (reverse) differentiation is itself an analytic approach at the system level, often used in topology optimization or sensitivity analysis to avoid recomputing a full Jacobian. 
One can use an adjoint equation to get the gradient of a single objective with respect to many inputs very efficiently -- but one still needs partial derivatives (Jacobian blocks) inside that adjoint equation. 
%
\cite{oberbichler2021efficient} apply the adjoint method to the global IGA
system, but they obtain the necessary
element-level derivatives via algorithmic differentiation of
the energy functional. 
The result was a clean separation: the
CAD-based geometry routines were automatically differentiated, and
then the adjoint method assembled the total sensitivity with far fewer
operations than a pure forward differentiation would require. This
hybrid of \autodiff{} and adjoint analytic methods allowed them to tackle shape
sensitivities for very high-order elements efficiently. 


% \subsection{Objectives}
% We aim to differentiate the FE element kernel
% analytically for efficiency, and let a higher-level \autodiff{} system handle
% glue code and complex couplings for flexibility.

% Pure AD-based FEM implementations exist and have shown what's possible (e.g.~end-to-end differentiable simulators like JAX-FEM and PyTorch-FEA). 
% However, they also highlight the downsides -- namely, the runtime and memory costs -- when applied to large, nonlinear FE problems. 
% On the other hand, decades of work on DDM in computational mechanics have furnished methods to get
% gradients efficiently for very specific problems (e.g.~the OpenSees DDM
% for structural sensitivities, or recent DDM for geometrically-nonlinear
% frame elements). 

% The hybrid approach seeks to get the ``best of
% both worlds.'' 
% By embedding an analytically differentiated finite element
% kernel inside a higher-level \autodiff{} framework, we retain the
% efficiency of hand-coded derivatives for the heavy lifting (the finite
% element internal calculations) and the flexibility of automatic
% differentiation for coupling with other software (optimization
% algorithms, machine learning models, etc.). 
% This approach has only started to appear in the literature in the last few years, and to our knowledge, no general-purpose Python-based finite element tool with built-in DDM
% sensitivities exists yet. 
% Our work will fill that gap. 
% We will demonstrate how a DDM-enhanced finite element solver can be embedded
% in an \autodiff{} environment to enable gradient-based design,
% optimization, and learning on complex FE models -- without the
% prohibitive memory cost of differentiating through the entire nonlinear
% solve. 
% This aligns with the trends seen in recent research and pushes
% them a step further by uniting the efficiency of classical methods with
% the flexibility of modern \autodiff{} frameworks.




